% ===================== MCMP117 Final Exam -- SAMPLE PAPER (student version) =====================
% Built with the course preamble:  python3 ../../lib/assemble_notes.py en 15 final_sample_body.tex MCMP117_Final_Sample.tex "Final Exam Sample Paper"
% Every answer is verified in ../code (theory_check.py, answers_sample.sh); see MCMP117_Final_AnswerKey.
\fancyhead[R]{\small\color{muted}Final Exam --- Sample Paper}
\setlength{\parskip}{0.35em}
\newcommand{\pts}[1]{\hfill\textcolor{mintdeep}{\textbf{[#1~pt]}}}
\newcommand{\ans}[1]{\par\smallskip\noindent\begin{tikzpicture}\draw[muted!60,rounded corners=2pt] (0,0) rectangle (\linewidth,#1);\end{tikzpicture}\par\medskip}
% empty time axis for a Gantt chart, 0..#1 ms
\newcommand{\ganttgrid}[2]{\par\noindent\begin{tikzpicture}[x=0.7cm]
\node[anchor=east,font=\small\bfseries] at (-0.2,0.31) {#2};
\draw[muted!70] (0,0) rectangle (#1,0.62);
\foreach \t in {0,...,#1} {\draw[muted!70] (\t,0) -- (\t,-0.1) node[below,font=\scriptsize,text=muted,inner sep=1pt] {\t};}
\end{tikzpicture}\par\smallskip}
% empty frame table: 12 references, 3 frames
\newcommand{\frametable}[1]{\begin{tabular}{@{}l|*{12}{>{\centering\arraybackslash}p{0.6cm}|}@{}}
\multicolumn{13}{@{}l}{\textbf{#1}}\\ \hline
ref & 5 & 4 & 1 & 5 & 0 & 5 & 4 & 0 & 2 & 1 & 0 & 1 \\ \hline
fr.\ 0 &&&&&&&&&&&& \\ \hline
fr.\ 1 &&&&&&&&&&&& \\ \hline
fr.\ 2 &&&&&&&&&&&& \\ \hline
fault &&&&&&&&&&&& \\ \hline
\end{tabular}}

\thispagestyle{empty}
\begin{tcolorbox}[enhanced,colback=dark,colframe=dark,arc=0pt,boxrule=0pt,
  left=18pt,right=18pt,top=14pt,bottom=14pt,coltext=white]
{\small\color{mint}\textsc{MCMP117 \,·\, Fundamentals of Operating Systems}}\par\smallskip
{\headfont\fontsize{22}{26}\selectfont\bfseries Final Exam --- Sample Paper}\par\smallskip
{\small Istanbul Okan University \,·\, Vocational School \,·\, Electronical Communication Technology Programme}\par
{\small Instructor: Asst.\ Prof.\ Dr.\ Hasan Oğuz \quad Date: [to be announced] \quad Duration: [90] minutes}
\end{tcolorbox}

\begin{center}
\renewcommand{\arraystretch}{1.6}
\begin{tabularx}{\linewidth}{@{}lX lX@{}}
\textbf{Name, surname:} & \hrulefill & \textbf{Student ID:} & \hrulefill \\
\textbf{Signature:} & \hrulefill & \textbf{Seat / room:} & \hrulefill \\
\end{tabularx}

\vspace{0.4em}
\renewcommand{\arraystretch}{1.35}
\begin{tabular}{|l|c|c|c|c|c|c|c|}
\hline
\textbf{Part} & A & B & C & D & E & F & \textbf{Total} \\ \hline
\textbf{Points} & 20 & 12 & 36 & 14 & 9 & 9 & \textbf{100} \\ \hline
\textbf{Score} & & & & & & & \\[0.4em] \hline
\end{tabular}
\end{center}

\textbf{Rules.} Closed book [to be confirmed]. [A basic, non-programmable calculator is allowed --- to be confirmed.] Phones off. Write your answers in the boxes; show every step of a calculation (formula $\to$ numbers $\to$ result with unit). Scheduling conventions of the course: context-switch time $0$; ties broken by arrival time, then by name; a smaller priority number means a higher priority; in Round Robin a process arriving at the moment a quantum expires enters the queue before the preempted process. TLB: use $\text{EAT} = t_{\text{TLB}} + (2-\alpha)\,t_m$. Linux questions refer to the practice folder \texttt{\textasciitilde/lab} (created by \texttt{mcmp117\_lab\_setup.sh}); your current directory is \texttt{\textasciitilde/lab} unless stated otherwise.

% ---------------------------------------------------------------------
\section*{Part A --- Multiple choice \hfill (10 $\times$ 2 = 20 points)}
Circle \emph{one} answer.
\begin{enumerate}[leftmargin=1.8em,label=A\arabic*.,itemsep=0.35em]
  \item UNIX was created at Bell Labs in \quad (a) 1956 \quad (b) 1969 \quad (c) 1981 \quad (d) 1991
  \item In which OS structure does only a minimal kernel run in kernel mode, with most services in user space? \quad (a) monolithic \quad (b) simple \quad (c) microkernel \quad (d) modular
  \item Which of these is \emph{not} stored in a process control block? \quad (a) process state \quad (b) program counter \quad (c) list of open files \quad (d) the source code of the program
  \item The time from a process's arrival until it first gets the CPU is its \quad (a) turnaround time \quad (b) waiting time \quad (c) response time \quad (d) burst time
  \item Which semaphore operation can block the calling thread? \quad (a) \texttt{signal(S)} \quad (b) \texttt{wait(S)} \quad (c) both \quad (d) neither
  \item The addresses generated by the CPU while a program runs are \quad (a) physical \quad (b) logical \quad (c) disk-block numbers \quad (d) inode numbers
  \item Thrashing means \quad (a) very high CPU utilisation \quad (b) constant page faults and low CPU utilisation \quad (c) a TLB hit ratio of $100\,\%$ \quad (d) a full disk
  \item After \texttt{chmod 754 f}, \texttt{ls -l f} shows \quad (a) \texttt{-rwxr-xr--} \quad (b) \texttt{-rwxr-x---} \quad (c) \texttt{-rwxrw-r--} \quad (d) \texttt{-rw-r-xr--}
  \item \texttt{cmd 2> /dev/null} \quad (a) discards the normal output \quad (b) discards the error messages \quad (c) writes errors to a file named \texttt{2} \quad (d) appends the output
  \item Which command lists the background jobs of the current shell? \quad (a) \texttt{ps} \quad (b) \texttt{top} \quad (c) \texttt{jobs} \quad (d) \texttt{bg}
\end{enumerate}

% ---------------------------------------------------------------------
\section*{Part B --- Short answers \hfill (3 $\times$ 4 = 12 points)}
\begin{enumerate}[leftmargin=1.8em,label=B\arabic*.]
  \item A program calls \texttt{write()}. Describe, step by step, how the call reaches the kernel and returns (API/wrapper, system-call number, trap, mode bit, system-call table). \pts{4}
  \ans{3.2cm}
  \item What is a race condition? Use \texttt{counter++} as the example, and name the three requirements that a solution to the critical-section problem must satisfy. \pts{4}
  \ans{3.2cm}
  \item What is thrashing, and how does the working-set model help the OS to prevent it? \pts{4}
  \ans{3.2cm}
\end{enumerate}

% ---------------------------------------------------------------------
\newpage
\section*{Part C --- Numerical problems \hfill (36 points)}
\begin{enumerate}[leftmargin=1.8em,label=C\arabic*.]
  \item Four processes (times in ms):
  \begin{center}
  \begin{tabular}{lcccc}
  \toprule
  Process & P1 & P2 & P3 & P4 \\ \midrule
  Arrival & 0 & 2 & 3 & 5 \\
  Burst & 6 & 2 & 7 & 3 \\
  Priority & 3 & 1 & 4 & 2 \\ \bottomrule
  \end{tabular}
  \end{center}
  For each algorithm draw the Gantt chart and compute the average waiting time and the average turnaround time: (a) FCFS, (b) preemptive Priority, (c) Round Robin with $q = 2$. \pts{12}
  \ganttgrid{18}{FCFS}
  \ganttgrid{18}{Priority}
  \ganttgrid{18}{RR, $q=2$}
  \ans{4.6cm}
  \item Pages are $2$~KB. The page table is: page $0 \to$ frame $3$, $1 \to 8$, $2 \to 6$, $3 \to 1$ (four entries).
  (a) Translate the logical address $5000$ into a physical address. \pts{3}\\
  (b) What happens for the logical address $9000$? \pts{1}\\
  (c) $t_{\text{TLB}} = 10$~ns, $t_m = 90$~ns, $\alpha = 0.8$: compute the effective access time. \pts{2}\\
  (d) Free holes of $200$, $450$, $100$, $300$, $500$~KB (in this order). Which hole receives a request of $280$~KB with first-fit, best-fit and worst-fit? \pts{2}
  \ans{5.2cm}
  \Needspace{20\baselineskip}
  \item Reference string $5, 4, 1, 5, 0, 5, 4, 0, 2, 1, 0, 1$ with $3$ frames. Fill in the frame tables and give the number of page faults for FIFO, OPT and LRU. \pts{8}
  \begin{center}
  \small\renewcommand{\arraystretch}{1.2}
  \frametable{FIFO}\quad faults: \underline{\hspace{1.2cm}}\\[0.5em]
  \frametable{OPT}\quad faults: \underline{\hspace{1.2cm}}\\[0.5em]
  \frametable{LRU}\quad faults: \underline{\hspace{1.2cm}}
  \end{center}
  \item (a) $t_m = 100$~ns, $t_{\text{fault}} = 5$~ms, page-fault rate $p = 0.0005$: compute the effective access time. How many times slower than RAM is it? \pts{3}\\
  (b) A file has records of $50$~bytes and blocks of $512$~bytes, both numbered from $0$. In which block and at which offset does record $100$ start? Does the record fit in that block? \pts{3}\\
  (c) Jobs wait for I/O half of the time ($p = 0.5$): compute $U(4)$. A program is $25\,\%$ serial: what speedup does Amdahl's law allow on $3$ cores, and what is the limit for many cores? \pts{2}
  \ans{6.0cm}
\end{enumerate}

% ---------------------------------------------------------------------
\newpage
\section*{Part D --- Write the command \hfill (7 $\times$ 2 = 14 points)}
Write \emph{one} command line for each task.
\begin{enumerate}[leftmargin=1.8em,label=D\arabic*.]
  \item Long listing of \texttt{docs}, including hidden files, with human-readable sizes. \ans{0.9cm}
  \item Rename \texttt{docs/notes\_01.txt} to \texttt{docs/meeting\_01.txt}. \ans{0.9cm}
  \item Create the empty files \texttt{report1.txt} \dots\ \texttt{report5.txt} with one short command. \ans{0.9cm}
  \item Create a symbolic link \texttt{latest.log} in the current directory that points to \texttt{logs/access.log}. \ans{0.9cm}
  \item Add execute permission for the owner (only) of \texttt{scripts/backup.sh}, in symbolic mode. \ans{0.9cm}
  \item The process with PID $4321$ ignores SIGTERM. Stop it anyway. \ans{0.9cm}
  \item Show how much disk space the directory \texttt{logs} uses, in human-readable form. \ans{0.9cm}
\end{enumerate}

\section*{Part E --- Write the pipeline \hfill (3 $\times$ 3 = 9 points)}
\begin{enumerate}[leftmargin=1.8em,label=E\arabic*.]
  \item \texttt{data/cities.txt} has one city per line. Print the most frequent city together with its count. \ans{1.1cm}
  \item How many requests in \texttt{logs/access.log} were for the page \texttt{/login}? (A line looks like\\ \texttt{10.0.1.10 - - [12/Oct/2026:08:00:00 +0300] "GET /about.html HTTP/1.1" 404 442}.) \ans{1.1cm}
  \item \texttt{data/students.csv} has the header \texttt{id,name,department,midterm,final}. Print the name and final grade (\texttt{name,final}) of the three students with the highest final grades, best first. \ans{1.1cm}
\end{enumerate}

\section*{Part F --- Write a script \hfill (9 points)}
Write \texttt{passfail.sh}. It takes exactly one argument, a threshold. Without exactly one argument it prints a usage message to standard error and exits with status $1$. Otherwise it counts how many students in \texttt{data/students.csv} have a final grade greater than or equal to the threshold (``passed'') and how many have less (``failed''), and prints for example \texttt{threshold 50: 30 passed, 10 failed}.
\ans{9.0cm}

\begin{center}\itshape --- End of the paper. Check your answers. Good luck! ---\end{center}
